\chapter{Thermodynamics}
\label{chap:thermodynamics}

% Closed question: Why does the record only grow, and when is it closed enough
% to support inference?

\begin{chapteressay}

Drive one mile forward and one mile backward. The car returns to its starting
place. The odometer does not return to its starting value. It reads two more
miles.

That is thermodynamics in dashboard form. Some coordinates can reverse. A
record of physical expenditure does not simply undo itself. The tire warmed.
The engine burned fuel or the battery discharged. Bearings heated. Road noise
became heat. The odometer did not measure displacement. It measured accumulated
path.

The first metaphor is therefore append-only record. Entropy is not just
messiness as a mood. It is the growth of physically unavailable information:
the spreading of energy into degrees of freedom that cannot be gathered back
without leaving further marks. A log can burn. Smoke and heat can disperse.
The reverse movie is not forbidden because the frames cannot be imagined. It
is forbidden because the required record-gathering would cost more record.

The second example is a cash register at a bakery. A loaf is sold. The drawer
opens. The receipt prints. Flour has already been baked, heat has already
dissipated, the loaf has left the shelf. If the customer returns it, the sale
is not erased. A refund is entered. The ledger grows. Closing the register at
night does not undo the day; it reconciles the day well enough to support
inference.

The code names this with comic force. \lean{Bullshit} has zero, one, and rest:
origin, first burden, and accumulated burden. \lean{AtreyuProcess} satirizes a
measurement into a larger record. \lean{TrueOutput} compares outputs by strict
growth. \lean{Closure} then asks when two accumulated records can be declared
same, different, or inferred. \lean{INFERRED\_TRUE} is the closure gesture:
the theory becomes true enough when the record has been closed by its own
equivalence process.

The closed question is:

\begin{center}
\emph{Why does the record only grow, and when is it closed enough to support
inference?}
\end{center}

The answer is that erasure is physical. Records can be transformed, summarized,
refunded, compressed, or closed, but not removed for free. Inference begins
when the growing ledger has been reconciled into an equivalence strong enough
to act on.

Chapter 11 closes Volume 2 by making measurement irreversible. The physical
gauge is not merely a way to compare. It is a way to carry accumulated cost.

\end{chapteressay}

% Phenomenon arcs use: 1/3 setup -> phenomenon box -> 2/3 structural interpretation.

\section{Irreversibility}
\label{se:irreversibility}

The odometer is the clean instrument. It accumulates path length. It does not
subtract when the car reverses. A round trip has zero net displacement and
positive mileage.

\begin{phenom}{Append-Only Odometer Effect}
\label{ph:append-only-odometer-effect}

\PhStatement A physical record can grow monotonically even when the coordinate
state it describes returns to its starting value.

\PhOrigin Odometers and thermodynamic entropy both distinguish accumulated
process from reversible coordinate description.

\PhObservation Driving out and back returns the car to its starting position
while increasing the odometer reading and dissipating energy into heat.

\PhConstraint The measurement may not confuse reversible displacement with
reversible physical expenditure.

\PhConsequence Irreversibility is ledger growth: the path has been carried even
when the endpoint has returned.

\end{phenom}

This gives entropy its local grammar. The second law says that the total
entropy of an isolated system does not decrease. In this book's language, the
complete physical ledger does not shrink. Local order can increase, but only by
exporting burden to the surrounding record.

\section{Landauer's Principle and Erasure Cost}
\label{se:landauers-principle-and-erasure-cost}

Landauer's principle makes erasure physical. To erase one bit of information at
temperature $T$, the minimum dissipated heat is

\[
Q \geq k_B T \ln 2.
\label{eq:chapter11-landauer}
\]

The bit may look abstract on a blackboard. In a device, it is a physical state.
Resetting it compresses alternatives into one state, and the lost distinction
must go somewhere.

\begin{phenom}{Landauer Erasure Effect}
\label{ph:landauer-erasure-effect}

\PhStatement Erasing a record has a minimum thermodynamic cost because the
distinction removed from memory must be exported to the environment.

\PhOrigin Landauer identified the heat cost of logically irreversible erasure;
later experiments measured this cost in physical memory systems.

\PhObservation Resetting a one-bit memory to a standard state dissipates at
least $k_B T \ln 2$ of heat in the ideal limit.

\PhConstraint The computation may not treat logical deletion as physically free
when a real device must implement the reset.

\PhConsequence Information is a physical ledger entry. Removing it from one
register writes cost into another.

\end{phenom}

This is the thermodynamic face of compilation. A compiler may remove dead code
from an output, but the process of computing, storing, comparing, and erasing
is carried by hardware. The visible artifact can be smaller while the physical
ledger still grows.

\section{Maxwell's Demon}
\label{se:maxwells-demon}

Maxwell's demon imagines a tiny observer sorting fast and slow molecules to
make one side of a box hotter than the other. It looks like entropy decreases
for free. The missing cost is the demon's memory.

To sort, the demon must measure. To repeat, the demon must remember or erase.
Bennett's resolution is that the thermodynamic bill arrives when the demon's
memory is reset.

\begin{phenom}{Maxwell Demon Effect}
\label{ph:maxwell-demon-effect}

\PhStatement A local entropy decrease obtained by measurement must include the
memory record and erasure cost of the measuring agent.

\PhOrigin Maxwell's thought experiment challenged the second law; later
information-theoretic analysis located the cost in measurement memory and
erasure.

\PhObservation Sorting molecules requires acquiring and managing information
about their states, and resetting that information dissipates heat.

\PhConstraint The apparatus may not omit the observer's memory from the
thermodynamic ledger.

\PhConsequence The second law survives because the complete record includes the
demon, not only the gas box.

\end{phenom}

The demon is a warning against bad boundaries. If the system is drawn too
small, entropy appears to be cheated. Draw the boundary around gas, demon,
memory, and heat bath, and the ledger grows again.

\section{Dark Energy as Cosmological Boundary Example}
\label{se:dark-energy}

Type Ia supernovae revealed that the expansion of the universe is accelerating.
In the simplest cosmological model, the cosmological constant acts as a dark
energy component with negative pressure. The expansion history is a record of
which directions the universe can no longer reverse: the comoving distance
between distant galaxies is a quantity the universe accumulates and does not
unwind.

\begin{phenom}{Dark Energy Effect}
\label{ph:dark-energy-effect}

\PhStatement Cosmological expansion is an accumulated record whose large-scale
history cannot be treated as a reversible local displacement.

\PhOrigin Supernova observations in the late twentieth century showed that the
universe's expansion is accelerating.

\PhObservation Distant type Ia supernovae appear dimmer than expected in a
decelerating universe, supporting a positive cosmological constant or related
dark-energy component.

\PhConstraint Cosmology may not close the expansion ledger using matter-only
deceleration when the distance-redshift record demands acceleration.

\PhConsequence The record-growth arrow extends beyond laboratory heat. It also
appears in the large-scale record of cosmic history.

\end{phenom}

This section is a cosmological boundary example, not a closure of the
heat-engine ledger. Dark energy does not behave like a working fluid;
the thermodynamic tie at this scale is the record-growth arrow --- the
universe's accumulated history refusing to unwind --- and not the
Clausius entropy of a heat engine. The link to the chapter's earlier
sections is structural: both the laboratory ledger and the cosmological
record share the same direction of refinement. The mechanism is
different in each case.

\section{The Meissner Effect}
\label{se:the-meissner-effect}

A superconductor cooled below its critical temperature expels magnetic flux
from its interior. The Meissner effect is not merely perfect conductivity. It
is a new ordered phase with a closed electromagnetic response. The closure here
is \emph{phase closure} or \emph{order-parameter closure}, not entropy
closure: the ordered state restricts which field configurations are
admissible inside the bulk. The thermodynamic arrow at this temperature
is the irreversible accumulation of the symmetry-broken phase, not the
heat-engine entropy of the surroundings. The two closures share the
record-growth discipline; they differ in what the closing operation
forbids.

\begin{phenom}{Meissner Closure Effect}
\label{ph:meissner-closure-effect}

\PhStatement A phase transition can close a record by excluding a class of
admissible field configurations from the interior.

\PhOrigin The Meissner effect showed that superconductors expel magnetic flux
when entering the superconducting phase.

\PhObservation Below the critical temperature, magnetic fields are expelled
from the bulk rather than simply persisting forever as they would in an ideal
zero-resistance conductor.

\PhConstraint The material may not be described only as resistance-free wire;
its phase imposes a new closure condition on the field record.

\PhConsequence Closure can be structural: the phase changes which field
histories are admissible inside it.

\end{phenom}

The final section returns to inference. A closed phase, a closed register, a
closed equivalence class, and a closed theory all have the same book-shape.
They do not say that nothing happened. They say the record has reached a state
from which certain distinctions are no longer admissible inside the boundary.

That is why \lean{Closure.le} and \lean{EquivalenceProcess.close?} are the
right code names for the ending. The ledger grows. Then a closure process says
which grown records can be treated as same, different, or inferred. The theory
does not become true by escaping physical cost. It becomes usable by accounting
for cost well enough.

\begin{bridge}

The chapter's question was why the record only grows, and when it is closed
enough to support inference.

The answer is physical erasure cost and disciplined closure. The odometer
accumulates path. Landauer gives erasure a heat bill. Maxwell's demon fails
unless its memory is counted. Dark energy keeps cosmological history open to an
accelerating expansion record. The Meissner effect shows a phase closing the
admissible field histories inside a material.

This closes Volume 2. Measurement has become a physical gauge: a calibrated,
witnessed, frame-corrected, transport-sensitive, invariant, irreversible record
from which inference can be drawn.

\end{bridge}

\begin{coda}{A Bakery Cash Register at Closing Time}

The bakery opens before sunrise. Dough has been mixed in the dark. Ovens heat.
Steam clouds the windows. The first customer buys a roll and coffee. The
register prints a receipt.

Nothing dramatic has happened, but the ledger has grown.

By noon, the drawer contains bills, coins, card slips, refund receipts, and
notes about special orders. A tray of croissants sold out. Three loaves were
set aside for pickup. One cake was returned because the name was misspelled.
The baker does not erase the sale. She enters a refund and a waste note. The
ledger grows again.

The bread cannot become flour. The heat cannot become unburned gas. The sold
roll cannot be unsold by pretending the receipt never printed. Every correction
is another transaction. Even the attempt to clean the record leaves a record.

This is Landauer in the bakery. Erasing a line from the register would not make
the day simpler. It would require a manager key, an audit entry, and a reason.
The old distinction is removed from one view and written into another. The
system protects itself by refusing free deletion.

The demon appears as the cashier who knows which pastries are moving quickly
and which are going stale. She can sort the trays, discount the slow ones, and
keep the display attractive. But her knowledge has to be recorded: waste,
markdown, transfer, donation. If she acts without memory, she cannot improve
the day. If she keeps memory, the ledger grows.

At closing, the register is reconciled. Reconciliation is not reversal. It is
closure. The drawer count, card batch, waste sheet, and sales totals are made
compatible enough that the owner can infer the day's result. Profit, loss,
inventory, and tomorrow's bake are all drawn from the closed record.

The bakery performs the final typeclasses. \lean{Bullshit} is accumulated
burden. \lean{AtreyuProcess} turns the next measurement into a larger comic
truth. \lean{TrueOutput} compares grown outputs. \lean{Closure} decides when
records are same, different, or inferred. \lean{INFERRED} is the owner trusting
the closed till enough to order flour.

The lights go off. The smell of bread remains. The ledger does not shrink.

\end{coda}
